\documentclass[11pt,a4paper]{article}

\usepackage[a4paper,margin=2.2cm]{geometry}
\usepackage[utf8]{inputenc}
\usepackage[T1]{fontenc}
\usepackage{booktabs}
\usepackage{xcolor}
\usepackage{hyperref}
\usepackage{amsmath}

\hypersetup{
    colorlinks=true,
    linkcolor=black,
    urlcolor=blue
}

\setlength{\parindent}{0pt}
\setlength{\parskip}{0.8em}
\setlength{\emergencystretch}{2em}

\begin{document}

\begin{center}
    {\Large Informationsvisualisierung -- Übung 11: Sugiyama-Methode}
\end{center}

\section*{Lauffähiger Code}

\href{https://github.com/yaniktfl/iv-repo}{https://github.com/yaniktfl/iv-repo},
Ordner \texttt{11/}.

\section*{Ausführung}

Anders als in den vorigen Übungen reicht \texttt{elm reactor} hier nicht:
\texttt{elm reactor} kompiliert nur die \texttt{.elm}-Datei und liefert dafür
eine eigene, generierte HTML-Seite aus -- unser \texttt{index.html} mit dem
Dagre-Script-Tag und dem Ports-Glue-Code wird dabei komplett ignoriert. Ohne
diesen JS-Code läuft \texttt{dagre.layout} nie, und es erscheinen nur die
Überschriften, aber keine Graphen.

Stattdessen \texttt{Main.elm} bauen und den Ordner mit einem einfachen
HTTP-Server ausliefern:

\begin{verbatim}
cd 11
elm make Main.elm --output=main.js
python3 -m http.server 8731
\end{verbatim}

Anschließend im Browser \texttt{http://localhost:8731/index.html} öffnen.

\section*{Zusammenspiel Elm / JavaScript (Ports)}

Elm kennt Dagre nicht direkt; die Layout-Berechnung läuft über zwei Ports,
wie im Elm-Guide-Kapitel zu Ports beschrieben und in
\texttt{madsflensted/elm-dagre-example} vorgeführt:

\begin{verbatim}
port graphs : (Value -> msg) -> Sub msg
port layout : String -> Cmd msg
\end{verbatim}

Elm schickt über \texttt{layout} eine JSON-Beschreibung des Graphen
(Knoten/Kanten inkl. geschätzter Breite/Höhe) an JavaScript. In
\texttt{index.html} liest das Dagre selbst ein, berechnet das Layout und
schickt das Ergebnis (jetzt mit \texttt{x}/\texttt{y}-Positionen und
Kanten-Punkten) über \texttt{graphs} zurück an Elm:

\begin{verbatim}
app.ports.layout.subscribe(function (data) {
  var parsed = JSON.parse(data);
  var g = dagre.graphlib.json.read(parsed);
  dagre.layout(g);
  var graph = dagre.graphlib.json.write(g);
  graph.key = parsed.key;
  app.ports.graphs.send(graph);
});
\end{verbatim}

Da mehrere Layout-Anfragen gleichzeitig unterwegs sein können (Satz-Graph
plus drei Ranker-Varianten für Aufgabe 11.3), bekommt jede Anfrage ein
\texttt{key}-Feld, das JavaScript unverändert zurückspiegelt. Elm ordnet die
Antwort darüber wieder der richtigen Anfrage zu (\texttt{sentenceGraph} bzw.
\texttt{Dict String GraphData} für die Ranker-Varianten).

Ausgangspunkt war das erweiterte Minimal-Beispielprogramm (TypedSvg statt
Svg) aus der Vorlesung; die Knotengröße wird nicht mehr fest vorgegeben,
sondern vor dem Senden an Dagre aus der Label-Länge geschätzt
(\texttt{estimateNodeSize}), damit z.\,B. \texttt{"sentence"} eine breitere
Box bekommt als \texttt{"is"}.

\section*{Aufgabe 11.1: Sentence Tokenization (ohne CSS-Klassen)}

Die Kantenrichtung wird über einen SVG-\texttt{marker} sichtbar gemacht: Ein
\texttt{defs}-Block definiert einen Pfeilkopf, der jeder \texttt{polyline}
per \texttt{markerEnd "url(\#arrowhead)"} angehängt wird.

\begin{verbatim}
TypedSvg.marker
    [ SA.id "arrowhead"
    , SA.markerWidth (Px 8)
    , SA.markerHeight (Px 8)
    , SA.refX "6"
    , SA.refY "3"
    , SA.orient "auto"
    ]
    [ TypedSvg.path [ SA.d "M0,0 L6,3 L0,6 Z", SA.fill (Paint Color.black) ] [] ]
\end{verbatim}

CSS-Klassen der Knoten werden hier laut Aufgabenstellung noch nicht
berücksichtigt -- alle Knoten werden einheitlich (weißer Hintergrund, blauer
Rahmen) gezeichnet.

\section*{Aufgabe 11.2: Sentence Tokenization (mit CSS-Klassen)}

Derselbe Graph wird ein zweites Mal gezeichnet, jetzt mit den Original-CSS-Klassen
aus der Demo (\texttt{type-TOP}, \texttt{type-S}, \ldots, \texttt{type-TK} für
die Token-Blätter). Jeder Knoten trägt seine Klasse als \texttt{class}-Attribut
auf der \texttt{g}-Gruppe; ein per \texttt{TypedSvg.style} eingebettetes
\texttt{<style>}-Element im SVG übernimmt die Original-Regel der Demo:

\begin{verbatim}
g.type-TK > rect { fill: #00ffd0; }
\end{verbatim}

Dadurch werden die Token-Blätter ("This", "is", "an", "example",
"sentence") türkis eingefärbt, alle anderen Knoten bleiben weiß --
exakt wie im Original-Dagre-D3-Demo, nur dass Elm über CSS-Klassen und nicht
über eine eigene Fallunterscheidung im Elm-Code färbt.

\section*{Aufgabe 11.3: Folie 456 -- drei Ranker-Varianten}

Für diese Aufgabe wird der Graph von Folie~456 (15 Knoten, 0--14, 26
gerichtete Kanten) dreimal mit unterschiedlicher Dagre-\texttt{ranker}-Konfiguration
(\texttt{network-simplex}, \texttt{tight-tree}, \texttt{longest-path}) im
\texttt{value}-Feld der Layout-Anfrage an Dagre geschickt und untereinander
ausgegeben:

\begin{verbatim}
jsonGraphValues rankdir ranker =
    object
        [ ...
        , ( "ranker", Encode.string ranker )
        ]
\end{verbatim}

Der Graph wird dabei unverändert (im Ausgangszustand von Folie~456, vor der
händischen Kreisentfernung/Schichtung, die auf der Folie nur zur Illustration
gezeigt wird) an Dagre übergeben -- Kreisentfernung, Rangzuweisung, Ordnen
innerhalb der Ränge und Koordinatenzuweisung übernimmt \texttt{dagre.layout}
vollständig selbst. Genau das macht den Unterschied zwischen den drei
Ranker-Varianten sichtbar: Knoten~12 (eine Senke ohne ausgehende Kanten) landet
bei \texttt{network-simplex} und \texttt{tight-tree} in einem mittleren Rang,
bei \texttt{longest-path} dagegen ganz unten, weil dieser Ranker jeden Knoten
so spät wie möglich einordnet.

\end{document}
